\documentclass[11pt]{article}

% Nature Methods style packages
\usepackage[utf8]{inputenc}
\usepackage[T1]{fontenc}
\usepackage{amsmath,amssymb}
\usepackage{graphicx}
\usepackage{booktabs}
\usepackage{hyperref}
\usepackage{natbib}
\usepackage{xcolor}
\usepackage{siunitx}
\usepackage{lineno}
\usepackage{multirow}
\usepackage{array}
\usepackage{algorithm}
\usepackage{algorithmic}
\linenumbers

\hypersetup{
  colorlinks=true,
  linkcolor=blue!60!black,
  citecolor=blue!60!black,
  urlcolor=blue!60!black
}

\title{Lab Manager: Open-Source AI-Native Laboratory Operations \\
with Multi-Model Consensus Document Processing}

\author{
  Cong Shen\textsuperscript{1,2,*} \\[6pt]
  \textsuperscript{1}Shen Lab, Massachusetts General Hospital, Boston, MA, USA \\
  \textsuperscript{2}Harvard Medical School, Boston, MA, USA \\
  \textsuperscript{*}Corresponding author
}

\date{}

\begin{document}

\maketitle

% ============================================================
\begin{abstract}
Research laboratories generate large volumes of paper
documents---packing lists, invoices, certificates of analysis, and shipping
labels---containing critical supply chain, inventory, and regulatory
information. Manually transcribing these documents is error-prone and slow,
creating a bottleneck between receiving supplies and maintaining accurate
records. We present Lab Manager, an open-source, AI-native laboratory
operations platform that converts scanned documents into structured
database records via multi-model consensus.

The system employs a six-stage architecture: (0)~optical character
recognition, (1)~parallel extraction by three frontier vision-language
models (VLMs), (2)~field-level consensus voting with automatic conflict
resolution, (3)~cross-model review where each model validates the merged
extraction, (4)~rule-based validation, and (5)~human review for unresolved
conflicts. In a benchmark of 10 real documents from 10 vendors, individual
VLMs achieved 80--90\% accuracy on vendor names and purchase order numbers,
and the 2/3 majority consensus rule resolved all vendor name disagreements
without human intervention.

Deployed as a self-hosted Docker application, Lab Manager provides a
complete laboratory operations stack: inventory management, order tracking,
vendor databases, full-text search, analytics, role-based access control,
and audit trails. The system runs in production at Harvard Medical
School/Massachusetts General Hospital and is freely available under the
Apache-2.0 license at \url{https://github.com/labclaw/lab-manager}.
\end{abstract}

% ============================================================
\section{Introduction}
\label{sec:intro}

Research laboratories are document-intensive. A typical academic
lab receives dozens of supply shipments per week, each accompanied by packing
lists, invoices, certificates of analysis (COAs), or shipping labels that
encode vendor information, purchase order numbers, catalog numbers, lot
numbers, quantities, and pricing. Researchers need this information for
inventory management, regulatory compliance, grant accounting, and
reproducibility~\citep{wilkinson2016fair, baker2016reproducibility}. Yet
extracting it from paper documents remains a largely manual process.

Most laboratories still process supply documents by hand: a lab member reads
the packing list, transcribes key fields into a spreadsheet or notebook, and
files the paper. This workflow is slow, error-prone, and creates information
silos. Meanwhile, the broader trend toward autonomous and self-driving
laboratories~\citep{abolhasani2023rise, macleod2020self,
szymanski2023autonomous, adam2024automated, canty2025science} has
focused on experiment execution and synthesis, with less attention to the
upstream supply chain documentation that supports those experiments.
Laboratory information management systems
(LIMS)~\citep{trigg2008current} and electronic lab
notebooks~\citep{avery2000electronic, barillari2016openBIS,
schmid2016scinote} address some aspects of lab data management, but they
still require manual data entry and do not automate the
document-to-database conversion step. Recent reviews of total laboratory
automation~\citep{nam2025revolutionizing} similarly emphasize the need
for AI-driven data capture at the pre-analytical stage.

Vision-language models (VLMs) offer a potential solution by extracting
structured information directly from document images. Recent document AI
systems such as Nougat~\citep{blecher2023nougat},
DocLLM~\citep{wang2024docllm}, and
LayoutLMv3~\citep{huang2022layoutlmv3} have demonstrated strong
performance on academic and business document understanding, while
benchmarks such as
Image2Struct~\citep{roberts2024image2struct} and
LongDocURL~\citep{deng2025longdocurl} have revealed that even
frontier models like GPT-4o~\citep{achiam2023gpt4} and
Gemini~\citep{team2024gemini} leave significant gaps in structured
extraction accuracy. Claude~\citep{anthropic2024claude} and similar
models further advance these capabilities. However, no single model is
reliable enough for laboratory data, where errors in lot numbers,
quantities, or vendor identities can have downstream consequences for
regulatory compliance and experimental reproducibility. Prior work on
multi-model aggregation~\citep{wang2023self, jiang2023llm,
wang2024mixture} demonstrated that combining multiple model outputs
improves reliability, but these approaches have not been applied to
laboratory document processing.

Here we present Lab Manager, an open-source platform that addresses the
document-to-database bottleneck with a multi-model consensus pipeline. The
system runs three frontier VLMs in parallel on each scanned document,
applies field-level consensus voting to resolve disagreements, performs
cross-model review for additional error correction, and routes unresolved
conflicts to human reviewers. This architecture operationalizes the
principle that laboratory data requires human-in-the-loop
verification~\citep{monarch2021human, amershi2014power}---AI assists
extraction, humans confirm before data enters the database.

Our contributions are: (1)~a multi-model consensus architecture for
laboratory document extraction with three-tier conflict resolution
(unanimous, majority, human review); (2)~an open-source, full-featured
laboratory operations platform with 125 API endpoints covering inventory,
orders, vendors, search, analytics, and audit trails; (3)~a benchmark on
10 real laboratory documents from 10 vendors evaluating 4 VLMs on
field-level extraction accuracy; and (4)~a production deployment at
Harvard Medical School/Massachusetts General Hospital.

% ============================================================
\section{Results}
\label{sec:results}

\subsection{System overview}
\label{sec:overview}

Lab Manager is a web-based application for laboratory supply chain and
inventory management (Fig.~\ref{fig:overview}). The core workflow is:
(1)~a lab member uploads a scanned document image (packing list, invoice,
COA, or shipping label); (2)~the AI pipeline extracts structured data;
(3)~low-confidence fields enter a human review queue; and (4)~the system
stores approved records in the database and propagates them to inventory,
orders, analytics, and search indices.

The system is built on a three-tier architecture: a PostgreSQL 17 database
with 23 relational models, a FastAPI~\citep{fastapi2024} backend with 125
REST API endpoints across 21 route modules, and a React frontend. The
application provides Meilisearch for full-text search, role-based access
control (RBAC) with 7 permission levels, a natural-language query interface
(NL$\to$SQL via LiteLLM), CSV export, and a complete audit trail. Users
deploy it as a single Docker Compose stack.

\subsection{Multi-model consensus pipeline}
\label{sec:pipeline}

The document intake pipeline processes each uploaded image through six
stages (Fig.~\ref{fig:pipeline}):

\textbf{Stage 0: OCR.} Optical character recognition extracts raw text
from the document image. While modern document AI systems such as
Donut~\citep{kim2022donut} have explored OCR-free approaches, we use
OCR to provide a text fallback for the VLM extraction stage. The OCR
subsystem supports a tiered strategy:
local models (DeepSeek-OCR-2, PaddleOCR) for cost-free processing, cloud
APIs (Gemini Flash, NVIDIA NIM) for higher accuracy, and an automatic
fallback mode that attempts local models first. In a baseline evaluation
on 8 real laboratory documents, Apple Vision OCR achieved 93.8\% field-level
accuracy (61/65 fields) compared to 63\% for Tesseract (41/65).
Handwritten receiving dates remained the primary failure mode for both
systems.

\textbf{Stage 1: Parallel VLM extraction.} Three frontier VLMs process
the document image independently and concurrently using a
\texttt{ThreadPoolExecutor} with a 180-second timeout per model. Each
model receives the document image and a structured extraction prompt
specifying the target JSON schema (Table~\ref{tab:schema}). The prompt
encodes domain-specific rules developed from error analysis of real
documents---for example, that VCAT codes should not be mistaken for lot
numbers, that ``1,000'' on Bio-Rad forms denotes a decimal quantity of
1.000 (not one thousand), and that the vendor name is the supplier company,
not the shipping carrier.

\textbf{Stage 2: Consensus merge.} Field-level consensus voting resolves
inter-model disagreements through a three-tier strategy:
\begin{itemize}
  \item \textbf{Unanimous (3/3):} All models agree on a field value.
    The field is auto-accepted with high confidence.
  \item \textbf{Majority (2/3):} Two or more models agree. The majority
    value is accepted, and the dissenting model's value is recorded for
    audit purposes.
  \item \textbf{No consensus:} All models disagree. The field is set
    using a priority ranking (Claude Opus $>$ Gemini 3.1 Pro $>$ GPT-5.4)
    and flagged for human review.
\end{itemize}

\textbf{Stage 3: Cross-model review.} Each VLM receives the merged
extraction alongside the OCR text and is asked to verify each field.
Corrections proposed by two or more models are automatically applied.
This stage catches extraction errors that survived the consensus merge,
particularly for fields where the original image was ambiguous but the
OCR text is unambiguous.

\textbf{Stage 4: Rule-based validation.} A deterministic validator checks
the merged extraction against domain rules: vendor names must not contain
address keywords (``blvd,'' ``street,'' ``suite''); document types must
be one of 8 enumerated values; quantities must be positive and below
10,000; lot numbers must not match VCAT code patterns; and dates must be
valid ISO-8601 within a reasonable year range. Validation issues are
classified as \texttt{critical} (blocking) or \texttt{warning}
(informational).

\textbf{Stage 5: Human review.} Documents with unresolved conflicts
(\texttt{\_needs\_human = true}) or critical validation issues enter a
review queue. Human reviewers see the original document image, OCR text,
extracted fields, consensus metadata (which models agreed or disagreed on
each field), and validation warnings. Approved documents propagate to the
database; rejected documents are flagged for re-processing or manual entry.

\subsection{Iterative refinement}
\label{sec:refinement}

When the initial extraction confidence falls below 0.7 or validation
identifies critical issues, the pipeline enters an iterative refinement
loop. The VLM receives its own prior extraction alongside specific
feedback (e.g., ``Field `vendor\_name': looks\_like\_address'') and
produces a corrected extraction. Up to 2 refinement rounds are applied.
Refinement is logged with per-round confidence deltas, enabling audit
of how and why the extraction improved.

\subsection{VLM benchmark on real laboratory documents}
\label{sec:benchmark}

We evaluated 4 vision-language models on a benchmark of 10 real scanned
laboratory documents from 10 different vendors
(Table~\ref{tab:benchmark}). The documents include packing lists,
invoices, and shipping labels with varying layouts, print quality, and
content density. Ground truth was established by manual expert annotation.

Three key fields were evaluated: vendor name, purchase order (PO) number,
and document type.

\begin{table}[ht]
\centering
\caption{VLM extraction accuracy on 10 real laboratory documents from 10
  vendors. Accuracy is the fraction of documents with correctly extracted
  field values (case-insensitive, substring matching for vendor names).
  95\% Clopper--Pearson confidence intervals are shown in brackets.}
\label{tab:benchmark}
\small
\begin{tabular}{lcccc}
\toprule
\textbf{Model} & \textbf{Vendor} & \textbf{PO \#} &
  \textbf{Doc Type} & \textbf{Avg Time} \\
\midrule
Llama-4-Maverick 17B & 9/10 [55, 100] & 9/10 [55, 100] & 8/10 [44, 97] & 7.3\,s \\
Mistral Small 119B & 9/10 [55, 100] & 9/10 [55, 100] & 8/10 [44, 97] & 4.7\,s \\
Llama-3.2 90B & 8/10 [44, 97] & 4/10 [12, 74] & 8/10 [44, 97] & 14.6\,s \\
Nemotron Nano 12B & 8/10 [44, 97] & 6/10 [26, 88] & 8/10 [44, 97] & 7.1\,s \\
\bottomrule
\end{tabular}
\end{table}

Llama-4-Maverick and Mistral Small achieved the highest accuracy on
vendor name and PO number extraction (90\% each), with document type
classification at 80\%---matching all other models on that field. The
largest performance gap appeared on PO number extraction, where Llama-3.2
achieved only 40\%---a field that requires precise alphanumeric string
matching and is particularly sensitive to OCR errors. Document type
misclassification (20\%) occurred when invoices were mislabeled as packing
lists or vice versa.

We note that the 95\% confidence intervals for these proportions are
wide due to the small sample size ($n = 10$). For example, 9/10 and 8/10
accuracies yield overlapping CIs ([55\%, 100\%] vs.\ [44\%, 97\%]),
meaning the observed differences between Llama-4-Maverick, Mistral Small,
and Nemotron Nano on vendor name extraction are not statistically
distinguishable at $\alpha = 0.05$. Only the Llama-3.2 PO number
performance (CI: [12\%, 74\%]) is credibly lower than the top performers
(CI: [55\%, 100\%]), as these intervals do not overlap. A larger
benchmark ($n > 50$) is needed to resolve pairwise differences.

All models successfully parsed the extraction prompt and returned valid
JSON for at least 9/10 documents (parse rate 90--100\%). Item-level
extraction (catalog numbers, quantities, lot numbers) was evaluated but
showed greater variance due to the complexity of multi-row table parsing;
detailed item-level results are provided in the Supplementary Materials.

The consensus mechanism exploits inter-model complementarity:
when one model fails on vendor name extraction but the other two succeed,
the 2/3 majority rule resolves the field without human intervention. In
our benchmark, all four models achieved $\geq$80\% on vendor names
(Table~\ref{tab:benchmark}), meaning that at most one model fails on any
given document. Under this regime, the 2/3 majority rule would resolve
all vendor name fields without human review. The cross-model review stage
provides an additional correction opportunity when the merged extraction
conflicts with the OCR text.

\subsection{System scale and deployment}
\label{sec:deployment}

Lab Manager comprises approximately 67,000 lines of code: 21,300
lines of Python backend (119 source files), 3,000 lines of TypeScript/React
frontend (11 component files in a separate repository), and 43,000 lines of
Python tests (134 test files). The test suite includes unit tests, BDD
(behavior-driven development) specifications, end-to-end pipeline tests,
and a release-gate script that validates the complete user flow from login
through document upload, review, and inventory management.

The database schema comprises 23 SQLModel tables covering documents,
vendors, products, orders (with line items and order requests), inventory
(with lifecycle operations: consume, transfer, adjust, dispose, open),
locations, equipment, staff, audit logs, alerts, notifications, safety
incidents, usage events, consumption logs, and digest templates. All
list endpoints support pagination, filtering, and sorting, returning
standardized \texttt{\{items, total, page, page\_size, pages\}} response
envelopes.

The 125 API endpoints span 21 route modules (Table~\ref{tab:endpoints}),
covering the complete laboratory operations lifecycle from document intake
through inventory tracking, analytics, and export. Role-based access control
enforces 7 permission levels (viewer, member, editor, reviewer, manager,
admin, owner), with route-level guards on all sensitive endpoints.

\begin{table}[ht]
\centering
\caption{API endpoint coverage by category.}
\label{tab:endpoints}
\small
\begin{tabular}{lrl}
\toprule
\textbf{Category} & \textbf{Count} & \textbf{Key Operations} \\
\midrule
Documents     & 8  & Upload, review, approve/reject, stats \\
Vendors       & 7  & CRUD, linked products \& orders \\
Products      & 9  & CRUD, linked inventory \& orders \\
Orders        & 11 & CRUD, line items, receive shipment \\
Order Requests & 7 & Request, approve, track status \\
Inventory     & 14 & CRUD, consume, transfer, adjust, low-stock, expiring \\
Search        & 2  & Full-text query, autocomplete suggestions \\
NL Query      & 4  & Natural language $\to$ SQL translation \\
Analytics     & 11 & Dashboard, spending, inventory value, top products \\
Export        & 9  & CSV: inventory, orders, products, vendors \\
Alerts        & 5  & Low-stock/expiry alerts, acknowledge, resolve \\
Audit         & 2  & Activity log, record history \\
Notifications & 6  & In-app notifications, preferences \\
Team          & 8  & Invite, assign roles, deactivate \\
Equipment     & 5  & CRUD, maintenance tracking \\
Device        & 4  & Heartbeat, CRUD, status \\
Digest        & 5  & Scheduled digests, templates, settings \\
Safety        & 3  & Incidents, rules, checks \\
Barcode       & 1  & Scan and lookup \\
Import        & 3  & Bulk data import \\
Telemetry     & 3  & Usage tracking, analytics \\
\midrule
\textbf{Total} & \textbf{125} & \\
\bottomrule
\end{tabular}
\end{table}

Shen Lab (Massachusetts General Hospital / Harvard Medical School)
deploys Lab Manager in production on a DigitalOcean droplet via Docker
Compose, using Caddy as the reverse proxy and an optional Cloudflare Tunnel
for secure access. A one-command installer (\texttt{deploy/install.sh})
enables deployment by non-technical users on Ubuntu or Debian systems.

% ============================================================
\section{Discussion}
\label{sec:discussion}

Lab Manager addresses a gap in the laboratory software landscape: the
conversion of paper documents into structured, queryable data. While LIMS
and ELN systems provide data management infrastructure, they assume data
arrives in structured form. Lab Manager automates the critical upstream
step of \emph{document understanding}, turning scanned packing lists and
invoices into database records that feed into inventory, ordering, and
analytics workflows.

\textbf{Multi-model consensus as a design pattern.} The core architectural
insight is that no single VLM is reliable enough for laboratory data, but
the \emph{agreement of multiple models} provides a useful confidence
signal. Individual models achieved 80--90\% accuracy on vendor names
(Table~\ref{tab:benchmark}); when two or three models independently agree on
a field value, the effective error rate drops further because the models
make different errors. When models disagree, the system flags the field
for human review rather than guessing. This pattern---parallel extraction,
consensus voting, cross-review, and human escalation---is generalizable
to any domain where AI-assisted data extraction requires high reliability.

\textbf{Domain-specific prompt engineering.} Accurate extraction from
laboratory documents requires extensive domain knowledge encoded in the
extraction prompt. We identified four recurring failure modes through iterative
error analysis on real documents: VCAT institutional codes mistaken for
lot numbers, shipping carrier names confused with vendor names, decimal
quantities misinterpreted as thousands, and address blocks extracted as
vendor identifiers. Each failure mode was addressed by adding explicit
rules to the extraction prompt and corresponding checks to the rule-based
validator.

\textbf{Limitations.} The current benchmark evaluates individual
VLMs on 10 documents but does not report end-to-end accuracy of the
consensus pipeline; a follow-up evaluation quantifying consensus
accuracy, human review rates, and error correction rates is needed.
The benchmark is limited to 10 documents from 10 vendors; a larger
evaluation across 100+ documents and 50+ vendors is needed to establish
statistical confidence in accuracy metrics. The consensus pipeline
requires 3 API calls per document, which increases cost and latency
compared to single-model extraction; the tiered OCR strategy and local
model options partially mitigate this. Handwritten text (particularly
receiving dates) remains the primary failure mode for OCR, with even
the best local model (Apple Vision OCR) achieving only 1/5 accuracy on
handwritten dates. The system sends laboratory documents to three cloud
VLM APIs, which raises data privacy concerns for labs handling
confidential reagents or proprietary materials; an on-premise VLM
deployment option would address this. The system has been tested with a
limited set of document types common in biomedical research
laboratories; generalization to chemistry, physics, or engineering
supply documents requires validation.

\textbf{Future directions.} We plan to expand the benchmark to 279
real documents from the production deployment and report field-level
precision, recall, and F1 scores. Multi-seed statistical analysis across
VLM providers will establish confidence intervals for accuracy claims.
Integration with PubChem for chemical product enrichment (CAS number
lookup, hazard information) is already implemented and will be evaluated
in a follow-up study. We are exploring fine-tuning smaller VLMs
(Nemotron Nano 12B, Llama-3.2 90B) on laboratory document extraction to
reduce dependence on frontier API models while maintaining accuracy.

% ============================================================
\section{Methods}
\label{sec:methods}

\subsection{Extraction schema}

The structured extraction schema (Table~\ref{tab:schema}) defines 14
document-level fields and 9 item-level fields, implemented as Pydantic
models with type validation and field-level constraints.

\begin{table}[ht]
\centering
\caption{Extraction schema for laboratory documents.}
\label{tab:schema}
\small
\begin{tabular}{lll}
\toprule
\textbf{Level} & \textbf{Field} & \textbf{Type} \\
\midrule
\multirow{14}{*}{Document}
  & vendor\_name & string (nullable) \\
  & document\_type & enum (8 types) \\
  & po\_number & string (nullable) \\
  & order\_number & string (nullable) \\
  & invoice\_number & string (nullable) \\
  & delivery\_number & string (nullable) \\
  & order\_date & ISO date (nullable) \\
  & ship\_date & ISO date (nullable) \\
  & received\_date & ISO date (nullable) \\
  & received\_by & string (nullable) \\
  & ship\_to\_address & string (nullable) \\
  & bill\_to\_address & string (nullable) \\
  & items & array of Item \\
  & confidence & float [0, 1] \\
\midrule
\multirow{9}{*}{Item}
  & catalog\_number & string (nullable) \\
  & description & string (nullable) \\
  & quantity & decimal (nullable) \\
  & unit & string (nullable) \\
  & lot\_number & string (nullable) \\
  & batch\_number & string (nullable) \\
  & cas\_number & string (nullable) \\
  & storage\_temp & string (nullable) \\
  & unit\_price & decimal (nullable) \\
\bottomrule
\end{tabular}
\end{table}

\subsection{Consensus algorithm}

The consensus merge operates on a dictionary of model extractions
$\{m_i: E_i\}_{i=1}^{N}$ where $N$ is the number of successful
extractions ($N \leq 3$, excluding models that timed out or failed).
For each field $f$:

\begin{enumerate}
  \item \textbf{Normalize}: Convert each value $E_i[f]$ to a canonical
    JSON string for comparison.
  \item \textbf{Count}: Group models by identical normalized values.
    Let $G_1, G_2, \ldots$ be the groups sorted by size (descending).
  \item \textbf{Resolve}:
    \begin{itemize}
      \item If $|G_1| = N$: unanimous. Accept value, mark
        \texttt{agreement: "unanimous"}.
      \item If $|G_1| \geq 2$ and $|G_1| > |G_2|$: majority. Accept
        $G_1$'s value, record dissenting values.
      \item If $|G_1| = |G_2|$ (tie with $\geq 2$ models each): flag
        \texttt{needs\_human: true}, use priority ranking as tiebreaker.
      \item Otherwise (all disagree): flag \texttt{needs\_human: true},
        select value from highest-priority model.
    \end{itemize}
\end{enumerate}

The priority ranking reflects empirically observed extraction quality:
Claude Opus 4.6 $>$ Gemini 3.1 Pro $>$ GPT-5.4. When all models fail
(all timeouts/errors), the entire document is flagged for manual
processing.

\subsection{Validation rules}

The rule-based validator implements 6 rule categories:
\begin{itemize}
  \item \textbf{Vendor name}: Must be $\leq$100 characters; must not
    contain address keywords (blvd, street, ave, suite, road, drive) or
    template text (``provider:'', ``organization'').
  \item \textbf{Document type}: Must be one of 8 enumerated values
    (packing\_list, invoice, certificate\_of\_analysis, shipping\_label,
    quote, receipt, mta, other).
  \item \textbf{Quantities}: Must be non-negative and non-zero
    (critical); warning above 10,000.
  \item \textbf{Lot numbers}: Must not match VCAT code patterns
    (critical).
  \item \textbf{Dates}: Must parse as valid ISO-8601; year must be
    within [current year $- 6$, current year $+ 2$].
  \item \textbf{Content hash deduplication}: SHA-256 content hash
    prevents processing the same physical document twice.
\end{itemize}

\subsection{Benchmark setup}

The VLM benchmark used 10 real scanned laboratory documents selected to
represent the diversity of the production document corpus: packing lists,
invoices, and shipping labels from 10 different vendors (Air-Tite, Addgene,
ALSTEM, A-M Systems, ABclonal, ACROS Organics, ALDON, ATCC, AVANTIK,
Agilent Technologies). Ground truth for vendor name, PO number, document
type, and item count was established by manual expert annotation.

Four VLMs were evaluated via the NVIDIA NIM API:
Llama-4-Maverick-17B-128E-Instruct,
Nemotron-Nano-12B-v2-VL,
Mistral-Small-4-119B-2603, and
Llama-3.2-90B-Vision-Instruct.
Each model received the same extraction prompt and document image. Accuracy
was computed as exact match (case-insensitive) for PO number and document
type, and substring match (case-insensitive) for vendor name to accommodate
minor formatting differences (e.g., ``Addgene'' vs.\ ``Addgene, Inc.'').

\subsection{Technology stack}

\textbf{Backend}: Python 3.12+, FastAPI, SQLModel, PostgreSQL 17,
Alembic (migrations), SQLAdmin (admin interface), Meilisearch (full-text
search), LiteLLM (NL$\to$SQL), Instructor (structured LLM extraction),
structlog (structured logging).
\textbf{Frontend}: React, TypeScript, Vite, Tailwind CSS, Radix UI.
\textbf{Infrastructure}: Docker Compose, Caddy (reverse proxy),
Cloudflare Tunnel (optional).
\textbf{Testing}: pytest, pytest-asyncio, pytest-bdd, testcontainers
(PostgreSQL), ruff (linting), mypy (type checking), bandit (security
scanning).
\textbf{CI/CD}: GitHub Actions with lint, typecheck, test (Python 3.12
and 3.13), frontend build, e2e gate, conventional commit enforcement,
and a merge-gate workflow requiring all checks to pass.

\subsection{Availability}

Lab Manager is available under the Apache-2.0 license at
\url{https://github.com/labclaw/lab-manager}. The quickstart:
\begin{verbatim}
  git clone https://github.com/labclaw/lab-manager
  cd lab-manager
  bash scripts/bootstrap_local_env.sh "My Lab"
  docker compose up -d --build
\end{verbatim}
Then open \texttt{http://localhost} and complete the setup wizard.
Documentation, deployment guides, and API specifications are included
in the repository.

% ============================================================
\section{Data Availability}
The VLM benchmark data (10 documents, 4 models, per-field accuracy) is
available in the repository at
\texttt{benchmarks/nim-10doc-benchmark-20260324.json}. The evaluation
framework (\texttt{benchmarks/extraction\_eval/}) enables reproducible
benchmarking with custom document sets and ground truth annotations.
Raw document images are not included due to institutional data policies
but can be reproduced by any laboratory using their own scanned documents.

% ============================================================
\section{Code Availability}
The complete source code is available at
\url{https://github.com/labclaw/lab-manager} under the Apache-2.0
license. The repository includes the backend, frontend, Docker
deployment configuration, test suite, benchmark framework, and
documentation.

% ============================================================
\section{Acknowledgments}
We thank the members of Shen Lab at Massachusetts General Hospital for
testing the system with real laboratory documents and providing feedback
on the extraction accuracy. We acknowledge Google (Gemini API), Anthropic
(Claude API), OpenAI (GPT API), Meta (Llama models), NVIDIA (NIM
platform), and Mistral AI for providing access to their foundation models.
We built this work on open-source tools including FastAPI,
SQLModel, PostgreSQL, Meilisearch, React, and the Python scientific
computing ecosystem.

% ============================================================
\section{Author Contributions}
C.S.\ designed the system architecture, implemented the core pipeline,
conducted the VLM benchmark, and wrote the manuscript.

% ============================================================
\section{Competing Interests}
The author declares no competing interests.

% ============================================================
\bibliographystyle{naturemag}
\bibliography{references}

% ============================================================
% FIGURES
% ============================================================

\clearpage

\begin{figure}[ht]
  \centering
  \fbox{\parbox{0.9\textwidth}{\centering\vspace{2.5cm}
    \textbf{Figure 1: Lab Manager system overview}\\[6pt]
    To be produced: system architecture diagram showing the four-layer stack:\\
    Document Upload $\to$ AI Pipeline $\to$ Human Review $\to$ Database\\[6pt]
    Shows: React frontend, FastAPI backend, PostgreSQL + Meilisearch,\\
    Docker Compose deployment. 125 API endpoints across 21 modules.
  \vspace{2.5cm}}}
  \caption{\textbf{Lab Manager system architecture.}
  The platform comprises a React frontend, FastAPI backend with 125 REST
  API endpoints, PostgreSQL database with 23 relational models, and
  Meilisearch full-text search. The system deploys as a single Docker
  Compose stack with Caddy reverse proxy. Users upload scanned documents
  through the web interface; the AI pipeline extracts structured data;
  human reviewers verify uncertain fields; approved records propagate to
  inventory, orders, analytics, and search indices.}
  \label{fig:overview}
\end{figure}

\begin{figure}[ht]
  \centering
  \fbox{\parbox{0.9\textwidth}{\centering\vspace{3cm}
    \textbf{Figure 2: Multi-model consensus document intake pipeline}\\[6pt]
    To be produced: 5-stage pipeline diagram:\\
    Stage 0 (OCR) $\to$ Stage 1 (Parallel VLM Extraction) $\to$\\
    Stage 2 (Consensus Merge: unanimous/majority/human) $\to$\\
    Stage 3 (Cross-Model Review) $\to$ Stage 4 (Validation) $\to$\\
    Stage 5 (Human Review Queue)\\[6pt]
    Shows: 3 VLMs running in parallel, consensus voting logic,\\
    iterative refinement loop, final human review gate.
  \vspace{3cm}}}
  \caption{\textbf{Multi-model consensus pipeline for document intake.}
  \textbf{Stage~0}: OCR extracts raw text (tiered: local or cloud).
  \textbf{Stage~1}: Three VLMs (Claude Opus, Gemini 3.1 Pro, GPT-5.4)
  extract structured data in parallel (180\,s timeout each).
  \textbf{Stage~2}: Field-level consensus voting---unanimous (3/3)
  auto-accepts, majority (2/3) auto-resolves, no consensus flags for
  human review. \textbf{Stage~3}: Cross-model review applies corrections
  when 2+ models agree on a change. \textbf{Stage~4}: Rule-based
  validation catches domain-specific error patterns (VCAT codes, address
  fragments, invalid dates). \textbf{Stage~5}: Human reviewers verify
  flagged fields with full context (original image, OCR text, consensus
  metadata, validation warnings).}
  \label{fig:pipeline}
\end{figure}

\begin{figure}[ht]
  \centering
  \fbox{\parbox{0.9\textwidth}{\centering\vspace{2cm}
    \textbf{Figure 3: VLM benchmark results}\\[6pt]
    To be produced: grouped bar chart showing per-model accuracy\\
    for vendor name, PO number, and document type on 10 documents.\\[6pt]
    Models: Llama-4-Maverick, Mistral Small, Llama-3.2, Nemotron Nano.\\
    Key finding: PO number extraction shows largest inter-model variance.
  \vspace{2cm}}}
  \caption{\textbf{VLM extraction accuracy on real laboratory documents.}
  Field-level accuracy of four VLMs on 10 real scanned documents from 10
  vendors. Vendor name and document type show relatively consistent
  performance across models (80--90\%), while purchase order number
  extraction varies widely (40--90\%), highlighting the value of
  multi-model consensus for this field. Error bars represent exact
  binomial 95\% confidence intervals.}
  \label{fig:benchmark}
\end{figure}

\begin{figure}[ht]
  \centering
  \fbox{\parbox{0.9\textwidth}{\centering\vspace{2cm}
    \textbf{Figure 4: Application screenshots}\\[6pt]
    To be produced: 4-panel screenshot montage:\\
    (a) Document upload and review queue\\
    (b) Inventory management with low-stock alerts\\
    (c) Analytics dashboard with spending trends\\
    (d) Mobile responsive view
  \vspace{2cm}}}
  \caption{\textbf{Lab Manager application interface.}
  \textbf{(a)}~Document review queue showing uploaded packing list with
  AI-extracted fields and consensus metadata. Fields with full agreement
  (green) are auto-accepted; fields requiring human verification (amber)
  are highlighted.
  \textbf{(b)}~Inventory management view with quantity tracking, expiry
  alerts, lot number traceability, and lifecycle operations (consume,
  transfer, adjust, dispose).
  \textbf{(c)}~Analytics dashboard showing spending trends by vendor,
  inventory value over time, and top-consumed products.
  \textbf{(d)}~Mobile-responsive layout for on-the-go inventory checks
  and barcode/QR scanning.}
  \label{fig:screenshots}
\end{figure}

\end{document}
